% !PRINTED 330
\weberpage{330}
%<<CONT
in which $h$, $a$, $b$, $c$ are constants, of which $h$ is different from zero,
and $f_3$ is a binary cubic form in $x_1, x_2$.

If, in accordance with this, we compute the function $H$ arranged
according to descending powers of $x_3$:

\weq{2}{H = x_3^{3n-6}\,H_0 + x_3^{3n-7}\,H_1 + \cdots,}

then we find

\weq{3}{\begin{aligned} H_0 &= -\,2\,(n-1)^2\,h^2\,c \\ H_1 &= -(n-1)^2\,h^2\,f_3''\,(x_2,x_2) + 4\,(n-1)(n-2)\,h\,(b^2-a\,c)\,x_1. \end{aligned}}

The curve represented by the equation $H=0$ is called
the \gesperrt{Hesse curve} of the curve $f$. The formulae (2) and (3)
show that the Hesse curve passes through the point $\xi_3$
when and only when $H_0=0$, that is $c=0$, i.\,e.\ by \S.\,86,
when this point is a point of inflection. From the expression
for $H_1$ one can conclude further that the inflectional tangent
$x_1=0$ is at the same time a tangent of the curve $H=0$
when and only when $f_3''\,(x_2,x_2)$, and consequently $f_3$ itself, is divisible
by $x_1$. In this case, as appears from (3), $x_1=0$ would be a
\gesperrt{tangent having four-point contact} at the point $\xi_3$. This
case can occur among curves of the third order only
when $f$ is divisible by $x_1$, hence never for an irreducible
curve of the third order.

Besides the points of inflection, the \gesperrt{Hesse} curve also passes
through the singular points of the curve $f$, which shall not
be investigated further here.

If, then, we restrict our considerations to curves $f$ without
a singular point, we can say that every point of
intersection of the curve $f$ with its \gesperrt{Hesse} curve gives a
point of inflection of the curve $f$, and that this curve $f$ has no
other points of inflection. If points with a tangent of four-point
contact occur, the curves $f$
and $H$ touch one another, and we can regard such points as
having arisen through the coincidence of two points of inflection. If
we further apply the theorem of Bezout (Bd.\,I, \S.\,49), we obtain
the result:

\begin{quote}
\gesperrt{A curve of the $n$th order without a singular point
has $3\,n\,(n-2)$ inflexion points;}
\end{quote}

and in particular:

\begin{quote}
\gesperrt{A curve of the third order without a singular point
has nine separately situated inflexion points.}
\end{quote}
